\documentclass[11pt,a4paper]{article}
\usepackage[margin=25mm]{geometry}
\usepackage{amsmath,amssymb,amsthm,hyperref}
\newtheorem{theorem}{Theorem}\newtheorem{lemma}[theorem]{Lemma}
\newcommand{\Aut}{\operatorname{Aut}}
\title{Simultaneous inversion and rank-three strings}
\author{Human Compression project}\date{}
\begin{document}\maketitle
\begin{abstract}
For a centreless nonsolvable group, the involution triple associated
with a reflexible generating pair satisfies the intersection property.
For an outer inverter, centrelessness alone suffices. We prove the
correspondences and specify how ordered orbit counts change when one
also identifies dual strings. This permits reuse of a string census
without assuming its counting convention agrees with a rotation table.
\end{abstract}

\section{The intersection property}
Let $G$ be centreless, and let $r,s\in G$ generate $G$, with
$|r|,|s|>2$ and $|rs|=2$. Suppose $b\in\Aut(G)$ inverts both
generators. Identify $G$ with its inner automorphisms. The accepted
inversion correspondence proves that $b$ is a unique nonidentity
involution, and that
\[
 a=rb,\qquad c=bs
\]
are nonidentity involutions with $ab=r$, $bc=s$ and $ac=rs$.
Put $H=\langle a,b,c\rangle=\langle G,b\rangle$.
In particular $a$ and $c$ commute and are distinct. A string C-group
triple means that, in addition,
\[
 \langle x_i:i\in I\rangle\cap\langle x_j:j\in J\rangle
   =\langle x_k:k\in I\cap J\rangle
 \quad(I,J\subseteq\{0,1,2\}),\qquad(x_0,x_1,x_2)=(a,b,c).
                                                        \tag{1}
\]

\begin{lemma}\label{solvable}
If $r,s$ satisfy $(rs)^2=(rs^{-1})^2=1$, then
$\langle r^2,s^2\rangle$ is abelian and normal in $\langle r,s\rangle$,
and the quotient is an elementary abelian group of order at most four.
In particular $\langle r,s\rangle$ is solvable.
\end{lemma}
\begin{proof}
The two relations give $rs=s^{-1}r^{-1}$ and $rs^{-1}=sr^{-1}$.
Consequently
\[
 rs^2r^{-1}=(rs)(sr^{-1})=(s^{-1}r^{-1})(rs^{-1})=s^{-2}.
\]
Interchanging $r$ and $s$ gives $sr^2s^{-1}=r^{-2}$: the first
relation is unchanged by cyclic permutation, and the inverse of the
second is $(sr^{-1})^2=1$. Thus each generator normalizes
$N=\langle r^2,s^2\rangle$. Also $r^2$ centralizes $s^2$, since
$r$ inverts $s^2$, so $N$ is abelian. Modulo $N$, the images of
$r$, $s$ and $rs$ are involutions or the identity. The images of
$r$ and $s$ therefore commute, proving the quotient assertion.
\end{proof}

\begin{theorem}\label{ip}
If $b$ is outer, or if $G$ is nonsolvable, then $(a,b,c)$ is a
string C-group triple in $H$.
\end{theorem}
\begin{proof}
First $\langle r\rangle\cap\langle s\rangle=1$. Indeed an element
of this intersection commutes with both $r$ and $s$, hence belongs
to $Z(G)=1$. The two adjacent dihedral subgroups are
\[
 D_0=\langle a,b\rangle=\langle r,b\rangle,
 \qquad D_1=\langle b,c\rangle=\langle s,b\rangle.
\]
Every element of the first is $r^i b^\epsilon$, and of the second
is $s^j b^\delta$, where $\epsilon,\delta\in\{0,1\}$.
These expressions are unique within each dihedral group: otherwise
$b$ would be a power of $r$ or $s$, would centralize that element,
and could not invert an element of order greater than two.

Suppose an element of $D_0\cap D_1$ lies outside $\langle b\rangle$.
If $\epsilon=\delta$, cancelling $b^\epsilon$ gives
$r^i=s^j=1$, a contradiction. If the parities differ, multiplying
the common element by $b$ if necessary gives
\[
                  z=r^i=s^j b.                         \tag{2}
\]
When $b$ is outer, (2) is impossible, since its left side belongs
to $G$ and its right side does not. In the other case $z$ is an
involution: it is a reflection in $\langle s,b\rangle$.
It centralizes $r$ and inverts $s$. Conjugating $(rs)^2=1$ by $z$
therefore gives $(rs^{-1})^2=1$. Lemma~\ref{solvable} would make
$G$ solvable, contrary to the hypothesis. Thus
$D_0\cap D_1=\langle b\rangle$ in either case.

For completeness this implies all of (1). The third pair subgroup
$V=\langle a,c\rangle$ is the four group. Neither $c$ nor $ac$
can belong to $D_0$: either would imply $c\in D_0\cap D_1$
and then $c=b$, contrary to $|bc|>2$. Hence
$V\cap D_0=\langle a\rangle$. Similarly
$V\cap D_1=\langle c\rangle$. The three involutions are distinct,
so their cyclic subgroups have trivial pairwise intersections.
Intersections involving the full set or the empty set are immediate.
These are all cases of (1).
\end{proof}

\section{The converse and duality}
\begin{theorem}\label{simple}
Let $G$ be nonabelian finite simple. Inner-reflexible generating pairs
$(r,s)$ of the above type correspond equivariantly to ordered rank-three
string C-group triples generating $G$. The correspondence is
$(r,s)\mapsto(rb,b,bs)$, with inverse $(a,b,c)\mapsto(ab,bc)$.
\end{theorem}
\begin{proof}
The forward assertion follows from Theorem~\ref{ip}, and
$H=G$ because $b\in G$. Conversely, let $(a,b,c)$ be such a
string triple, and put $r=ab$, $s=bc$. Then $rs=ac$ has order two.
The subgroup $K=\langle r,s\rangle$ is normalized by $b$, which
inverts both generators. Since $a=rb$ and $c=bs$, we have
$G=K\langle b\rangle$ and $[G:K]\leq2$. Here $K$ is nontrivial,
so simplicity gives $K=G$. The adjacent product orders exceed two:
otherwise an endpoint of the string would commute with both other
generators and be a nontrivial central element of $G$.
The unique-inverter statement gives the inverse correspondence,
and both directions commute with $\Aut(G)$.
\end{proof}
There is also an outer version. Fix a group $A$ with $G\lhd A$,
$[A:G]=2$, and faithful conjugation action on the centreless group
$G$. Reflexible generating pairs with inverter in $A\setminus G$
correspond to string C-group triples generating $A$ whose three
generators lie in $A\setminus G$. For the reverse direction the
rotation subgroup $K$ lies in $G$, and the same argument gives
$[A:K]\leq2$, hence $K=G$. The correspondence is equivariant for
the automorphisms of $A$ preserving $G$; in particular, for its
conjugation action. For $A=\Aut(G)$ this is the same action on
the pairs as in the accepted inversion bound.

Reversing a string sends $(a,b,c)$ to $(c,b,a)$. On rotations this is
\[
                   (r,s)\longmapsto(s^{-1},r^{-1}).       \tag{3}
\]
It is an involution commuting with automorphisms. Let $n_{p,q}$
count ordered automorphism orbits of the inner strings, with the
displayed orientation $p\leq q$. For $p=q$, let $d_p$ count those
orbits fixed by (3). If $Q$ denotes the number after also identifying
reversal, then
\[
 Q=\sum_{p<q}n_{p,q}+\sum_p\frac{n_{p,p}+d_p}{2}.        \tag{4}
\]
Indeed reversal exchanges the two different-order orientations,
whereas among equal-order orbits it has $d_p$ fixed points and
pairs all others. In particular the sum of the triangular rotation
table is
\[
 \sum_{p\leq q}n_{p,q}=Q+\sum_p\frac{n_{p,p}-d_p}{2},    \tag{5}
\]
not generally $Q$ or $2Q$. The count of all ordered strings, including
both different-order orientations, is $2Q-\sum_p d_p$.

The mixed-intersection obstruction and solvable-square argument were
proposed by the root agent and checked and developed here. The outer
case and the explicit reversal formulas complete the shared interface.
No new O'Nan census, exact rotation-table entry or source completion is
claimed by these theorems. They explain how any later accepted string
census may be reused, retaining its involution classes and fixed points
of reversal instead of identifying its total with a different count.
\end{document}
